% STATUS: verified
\chapter{Over time: the action register}\label{ch:register}

A compliance report describes one day. This chapter describes what the system keeps between
reports: the \term{action register}, the per-company list of the actions that the plans of its
signed-off reports called for, each with an owner, a due date, a state and the findings it is
to close. It states the problem and the one event that moves the register
(\cref{sec:register:why,sec:register:event}), what comes in at that event
(\cref{sec:register:plan}), the states of an action (\cref{sec:register:states}) and --- the
core of the chapter --- how later reports confirm and close actions
(\cref{sec:register:settle}). The identity of an obligation, on which this rests, the history
of a finding and the changes in the law follow
(\cref{sec:register:identity,sec:register:history,sec:register:law}); then what the chat
assistant and the Excel export do with the register (\cref{sec:register:reading}), where the
line between models, code and people runs (\cref{sec:register:line}), and what was tried on
real data (\cref{sec:register:tried}).

No model is called anywhere in this chapter. It therefore holds no path card: the register is
code and people, and \cref{sec:register:line} explains why that is deliberate.


\section{The problem: a plan that died with its report}\label{sec:register:why}

A report ends with an \term{action plan}: a prioritised list of actions, each naming the
findings it is to close, a timeframe and an owner role (\cref{ch:report} describes how it is
made). The plan is stored with the report, in \code{report_actions}. The next analysis of the
same company writes a new report with a new plan, and nothing connects the two: nobody can say
whether an action of the old plan was carried out, by whom, by when, or whether the gap it
addressed is gone.

The course's findings-and-actions register answers this with one rule: an action is closed
only after a re-test. Saying that the work was done is not enough; someone has to look again
and find the requirement met. The system already has a way of looking again --- the next
report --- and a person who takes responsibility for a report --- the one who signs it off.
The register joins the two:

\begin{itemize}
  \item it keeps the plan of every signed-off report as rows of its own, per company, in the
    table \code{company_actions} (model \code{CompanyAction}), which outlive the report;
  \item it uses the next signed-off report as the re-test: an action is closed when later
    signed-off reports find every obligation under it met.
\end{itemize}

All of it is in one file of the API, \file{analysis/register.py}, with seven public functions:
\code{sign_off}, \code{withdraw}, \code{standing}, \code{state_of}, \code{history},
\code{changes} and \code{law_changes}. Five constants at its head fix the vocabulary of the
whole chapter.

\begin{table}[H]
\caption{The constants at the head of the register's file}\label{tab:register:constants}
\begin{tabularx}{\textwidth}{@{}L{22mm}L{58mm}Y@{}}
\toprule
\hd{Constant} & \hd{Value} & \hd{Used for} \\
\midrule
\code{DUE_DAYS} & \code{immediate} 14, \code{30_days} 30, \code{90_days} 90, \code{12_months} 365 &
  Days from the sign-off to an action's due date, by the plan's timeframe. \\
\code{OPEN} & \code{missing}, \code{partial}, \code{unknown} & Statuses of a finding that needs an action. \\
\code{SETTLED} & \code{met}, \code{not_applicable} & Statuses that confirm a finding. \\
\code{LIVE} & \code{open}, \code{in_progress}, \code{done} & States of an action that is not closed. \\
\code{HISTORY} & 4 & How many earlier assessments a finding's history goes back. \\
\bottomrule
\end{tabularx}
\end{table}

\code{OPEN} and \code{SETTLED} sort the five statuses of a finding into the two groups the
register reasons with. A finding that is \status{missing}, \status{partial} or
\status{unknown} is \emph{open}: it needs an action, and a later report that finds it so
confirms nothing. A finding that is \status{met} or \status{na} is \emph{settled}. Wherever
the register reads a status it reads the \emph{effective} one --- a reviewer's override when
there is one, the model's verdict otherwise (\code{_effective}). \code{LIVE} names the three
states of an action that is not closed.


\section{The one event that moves the register}\label{sec:register:event}

The register changes its confirmations on one event only: a person \term{signing a finished
report off}, or that signature being taken back. The sign-off is the route
\route{PATCH}{/api/reports/{id}} with the body \code{{approved: true, note?}}; the console
sends it from the button \ui{Patvirtinti ataskaitą} (sign the report off) on the report's
page, whose dialog says, before the button, what the sign-off will do to the register
(\file{pages/Reports/ReportApproval.jsx}). Only a finished report can be signed off (409
\code{NOT_FINISHED} otherwise). The view \code{_approve} in \file{analysis/api/views.py} stamps
\code{approved_by}, \code{approved_at} and \code{approval_note} on the report and then calls
the register; the comment on that line reads \enquote{The company's register moves on a
sign-off and on nothing else}.

\begin{figure}[H]
\centering
\begin{tikzpicture}[node distance=3.5mm and 10mm]
  \node[man, text width=41mm] (sign) {signs a finished report off};
  \node[man, text width=41mm, below=of sign] (unsign) {withdraws the sign-off};
  \node[man, text width=41mm, below=of unsign] (review) {changes a finding of a signed-off report};
  \node[man, text width=41mm, below=of review] (delete) {deletes a signed-off report};
  \coordinate (mid) at ($(unsign.east)!.5!(review.east)$);
  \node[chk, text width=37mm, right=10mm of mid] (settle) {\code{_settle}\\ every confirmation worked out again from all signed-off reports};
  \node[chk, text width=30mm, right=of settle, yshift=13mm] (plan) {\code{sign_off}, step~2\\ at a sign-off only:\\ the report's plan comes in};
  \node[store, inner xsep=2pt, below=9mm of plan] (table) {\code{company_actions}};
  \node[man, text width=41mm, below=7mm of delete] (edit) {edits one action: owner, due date, state, note};
  \draw[flow] (sign.east) -- ++(4mm,0) |- ([yshift=3mm]settle.west);
  \draw[flow] (unsign.east) -- ++(4mm,0) |- ([yshift=1mm]settle.west);
  \draw[flow] (review.east) -- ++(4mm,0) |- ([yshift=-1mm]settle.west);
  \draw[flow] (delete.east) -- ++(4mm,0) |- ([yshift=-3mm]settle.west);
  \draw[flow] (settle.east) -- ++(4mm,0) |- (plan.west);
  \draw[flow] (settle.east) -- ++(4mm,0) |- (table.west);
  \draw[flow] (plan.south) -- (table.north);
  \draw[flow] (edit.east) -| (table.south) node[lbl, pos=.25] {\code{PATCH}~~\code{/api/actions/{id}}};
\end{tikzpicture}
\actorlegend
\caption{Everything that writes the register. Four acts of a person on a report reach the
one function that works the confirmations out; only a sign-off also brings a plan in. A
person may edit a single action directly. No job, no report being made or finished and no
uploaded document has a path into the table, and the chat has no tool that writes it: a
review it records on the user's instruction is the third act shown here
(\cref{sec:register:reading}).}\label{fig:register:events}
\end{figure}

\Cref{fig:register:events} shows every writer of the table. Two functions of
\file{analysis/register.py} are called from the views, and from nowhere else in the code:

\begin{table}[H]
\caption{The events that move the register, and what each calls}\label{tab:register:events}
\begin{tabularx}{\textwidth}{@{}L{48mm}L{37mm}Y@{}}
\toprule
\hd{Event} & \hd{Calls} & \hd{Effect} \\
\midrule
A person signs a finished report off (\code{_approve}, \code{approved: true}) &
  \code{register.sign_off} &
  The actions there are get held against all signed-off reports, this one now among them;
  then the report's plan comes in. Answers \code{{added, covered, confirmed, closed, reopened}}. \\
A person withdraws the sign-off (\code{_approve}, \code{approved: false}) &
  \code{register.withdraw} &
  The confirmations are worked out again without this report. Answers
  \code{{withdrawn}}, the number of actions that changed. \\
A reviewer changes a finding of a signed-off report (\code{finding_item}) &
  \code{register.withdraw} &
  Any change to a finding --- the overriding status, the reviewer's note, the cause --- clears
  the sign-off: a changed verdict needs a fresh one. What the sign-off confirmed goes back with it. \\
A signed-off report is deleted (\code{report_item}, \code{DELETE}) &
  \code{register.withdraw} &
  The sign-off is cleared first, the register is settled without the report, then the report
  is deleted: it takes with it what it alone confirmed. \\
\bottomrule
\end{tabularx}
\end{table}

The answer to a sign-off carries the report and, under \code{register}, the five counts; the
console words them in the success message with a link to the register, or says that the
register did not change when all five are zero. Both the sign-off and its withdrawal are
written to the audit log (\code{reports.approve}, \code{reports.unapprove}) with those counts.

\begin{humanbox}{Why only a signature moves it}
A finished report is a model's assessment that a person has not yet accepted. Its findings can
still be overridden in review, and the \code{Report} model describes the sign-off as a person
approving \enquote{the reviewed report as the company's assessment}. If an unsigned report
could close an action, a model's verdict would end a commitment that a person made. So the
evidence that closes an action is always a report somebody has put their name to, and when
that name is taken back --- by hand, or because a finding was changed afterwards --- the
closure goes with it. The register's own header says it shortly: a closure is a re-test, not
a statement.
\end{humanbox}

The same reasoning excludes the other candidates. A new report that has not been signed off
appears nowhere in the register; the test \code{test_a_sign_off_brings_the_plan_in_once}
asserts that an unsigned report \enquote{moves nothing}. No job handler imports the register's
writing functions. The chat assistant has one tool for the register and it only reads
(\cref{sec:register:reading}). Uploading a document that would close a gap changes nothing
until a later report has assessed it and a person has signed that report off.


\section{What comes in at a sign-off}\label{sec:register:plan}

\code{sign_off(report, user)} runs in one database transaction and in two steps. Step~1 holds
the existing actions against the signed-off reports (\cref{sec:register:settle}); step~2,
described here, brings the report's own plan in. It walks the plan's actions in their order
and, for each, builds the list of \term{items}: the findings the action is to close.

A finding of the plan action becomes an item unless one of three things is true:

\begin{enumerate}
  \item \textbf{It is not open.} Its effective status is not in \code{OPEN}: what is met needs
    no action.
  \item \textbf{A newer signed-off report already finds it settled.} \code{_later(report)}
    collects, for the company's signed-off reports made \emph{after} this one, the newest
    verdict on each obligation; when that verdict is in \code{SETTLED}, nothing is left to do.
    This matters when an older report is signed off late.
  \item \textbf{Its obligation is already under a live action.} The set \code{under_way}
    holds every obligation named by an item of an action in a \code{LIVE} state. Such a
    finding is counted as \code{covered} and left out: an obligation never stands under two
    live actions. An obligation added by this sign-off joins the set at once, so two actions
    of the same plan cannot both take it.
\end{enumerate}

An item is a snapshot, stored in the JSON column \code{items} of the action:
\code{obligationId}, \code{findingId}, \code{title} (the obligation's title as the finding
carries it), \code{act} and \code{path} (the act of the finding's first citation, or else the
finding's own act, and the provision of that citation), \code{sourceKind}, \code{status} (the effective status at the sign-off) and
\code{confirmed}, which starts as null. Like a report, the register copies the titles it
shows, so an action reads the same after the corpus or the report has changed.

A plan action whose findings were all left out is not brought in. Otherwise a
\code{CompanyAction} is created with the plan's \code{title}, \code{description},
\code{category}, \code{priority} and \code{effort}; the plan's owner role as \code{owner};
the state \code{open}; and a due date counted from the moment of the sign-off by the plan's
timeframe (\code{DUE_DAYS}): 14 days for \code{immediate}, 30 for \code{30_days}, 90 for
\code{90_days}, 365 for \code{12_months}, and 90 for a timeframe the table does not name. The due
date counts from the sign-off, not from the day the report was made, because the sign-off is
the day a person accepted the plan.

\paragraph{Nothing comes in twice.} Each register action remembers the plan action it came
from (\code{plan_action_id}), and a unique constraint on the company and that id backs the
check in the code. A sign-off withdrawn and given again therefore adds nothing; neither does
the next report, whose plan covers the same open obligations --- its findings are counted as
\code{covered}.

\paragraph{Actions that close no finding.} A plan may hold actions without findings ---
correcting the references the citation audit found wrong, finding out a fact the company's
profile does not state. They come in too, with an empty \code{items} list. Having no
obligation to be known by, they are matched by their \code{category} instead: while an
action without items from \emph{another} report is still live in a category, the next
report's action without findings in the same category is the same work asked again, is
counted as \code{covered} and is not brought in. No report can ever confirm such an action;
only a person closes it (\cref{sec:register:states}).


\section{The states of an action and what a person may change}\label{sec:register:states}

An action is stored in one of four states, \code{ACTION_STATES} in
\file{analysis/models.py}, which a database check constraint enforces.

\begin{table}[H]
\caption{The states of an action. The first four are stored; the fifth is worked out when the register is read}\label{tab:register:states}
\begin{tabularx}{\textwidth}{@{}L{22mm}L{45mm}Y@{}}
\toprule
\hd{State} & \hd{In the console} & \hd{Means} \\
\midrule
\code{open} & \ui{Atviras} (open) & Brought in by a sign-off; nothing reported yet. \\
\code{in_progress} & \ui{Vykdomas} (in progress) & Somebody is working on it. \\
\code{done} & \ui{Įvykdytas, laukia patikros} (done, awaiting re-test) &
  The owner's word that it was carried out. It is not a closure: the action waits for a later
  signed-off report. Setting it stamps \code{done_at}. \\
\code{closed} & \ui{Uždarytas} (closed) &
  Either later signed-off reports found every obligation under it met
  (\code{confirmed_report} names the closing report), or a person closed it by hand with a
  reason. \\
\midrule
\code{late} & \ui{Vėluojantis} (late) &
  Never stored. \code{state_of(action, today)} answers \code{late} for an action that is
  \code{open} or \code{in_progress} and whose due date is before today. \\
\bottomrule
\end{tabularx}
\end{table}

Late is derived because it depends on the calendar, not on anything that happened to the
action: a stored flag would be wrong the morning after the due date unless some job rewrote
it. An action called \code{done} is never shown as late; its owner has answered, and it waits
for the re-test. Every answer of the API carries both values: \code{status}, the stored
state, and \code{state}, the one to show.

\begin{figure}[H]
\centering
\begin{tikzpicture}[x=1mm, y=1mm]
  \node[man, text width=20mm] (in) at (11,0) {a report is signed off};
  \node[svc, minimum width=20mm] (open) at (47,0) {\code{open}};
  \node[svc, minimum width=26mm] (prog) at (84,0) {\code{in_progress}};
  \node[svc, minimum width=20mm] (done) at (121,0) {\code{done}};
  \begin{scope}[on background layer]
    \node[lane, fit=(open)(prog)(done), inner sep=7pt] (live) {};
  \end{scope}
  \node[draw=warn, dashed, rounded corners=3pt, line width=.6pt, fit=(open)(prog), inner sep=3.5pt] (late) {};
  \node[lanelabel, anchor=south west] at (live.north west) {the live states};
  \node[note, text=warn, anchor=north, text width=24mm, align=center] at (11,-8)
    {dashed frame: shown as \code{late} once the due date has passed};
  \draw[flow] (in) -- (open) node[lbl, pos=.4, above=1pt] {plan};
  \draw[flow, <->] (open) -- (prog) node[lbl, midway, above=1.5pt] {a person};
  \draw[flow, <->] (prog) -- (done) node[lbl, midway, above=1.5pt] {a person};
  \draw[flow] (done.north) -- ++(0,9) -| (prog.north)
    node[lbl, pos=.25, above=1pt] {a sign-off or a withdrawal: a report made after \code{done_at} still finds a finding open};
  \node[chk, text width=38mm] (crep) at (58,-31) {\code{closed} by a report\\ \code{confirmed_report} is set};
  \node[man, text width=38mm] (chand) at (116,-31) {\code{closed} by a person\\ a note says why};
  \coordinate (r1) at (52,0);  \coordinate (r2) at (64,0);
  \coordinate (h1) at (110,0); \coordinate (h2) at (122,0);
  \draw[flow] (r1 |- live.south) -- (r1 |- crep.north)
    node[lbl, midway, left=2pt, align=right] {a sign-off: every\\ finding confirmed};
  \draw[opt] (r2 |- crep.north) -- (r2 |- live.south)
    node[lbl, midway, right=2pt, align=left] {the evidence falls:\\ back to \code{closed_from}};
  \draw[flow] (h1 |- live.south) -- (h1 |- chand.north)
    node[lbl, midway, left=2pt, align=right] {a person,\\ with a note};
  \draw[flow] (h2 |- chand.north) -- (h2 |- live.south)
    node[lbl, midway, right=2pt, align=left] {a person\\ takes it back};
\end{tikzpicture}
\actorlegend
\caption{The states of an action and who moves it between them. A person moves an action
freely among the live states and may close it by hand; code closes it, sends it back to work
or reopens it only at a sign-off or its withdrawal. No arrow leads by hand out of a closure
made by a report.}\label{fig:register:states}
\end{figure}

\subsection{A person's hand on one action}

Between sign-offs a person keeps an action up through
\route{PATCH}{/api/actions/{id}} (\code{action_item}), sent by the form of an opened row on
the company's tab \ui{Veiksmai} (actions). The body carries only what changed.

\begin{table}[H]
\caption{What a person may change on an action}\label{tab:register:patch}
\begin{tabularx}{\textwidth}{@{}L{20mm}Y@{}}
\toprule
\hd{Field} & \hd{Rule} \\
\midrule
\code{owner} & Text, trimmed, at most 300 characters. It starts as the plan's owner role; a
  person replaces it with whoever really owns the work. \\
\code{dueDate} & A date, \code{YYYY-MM-DD}. \\
\code{note} & Text or null, trimmed, at most 4,000 characters. \\
\code{status} & One of \code{open}, \code{in_progress}, \code{done}, \code{closed} --- never
  \code{late}. A change to \code{done} stamps \code{done_at}. Closing needs a note. Reopening
  works only for a closure made by hand. \\
\bottomrule
\end{tabularx}
\end{table}

A wrong value is refused with 400 \code{INVALID_FIELD}; a user who cannot see the action's
company gets 404. Every change is written to the audit log as \code{actions.update} with the
fields that changed (the note's text is not copied there, only the fact that it changed).
There is no route that creates an action and none that deletes one: actions come from plans,
and they leave only with their company.

\begin{humanbox}{Closing by hand, and what a hand may not undo}
A person may close an action without any report's confirmation --- the service was
discontinued, the obligation no longer concerns the company --- but only with a note: without
one the API answers 400 \code{NOTE_REQUIRED} (\enquote{Say in the note why the action is
closed without a report's confirmation}), and the console keeps the save button off until
the note is written. The action then records who closed it and when (\code{closed_by},
\code{closed_at}) and the state it was closed from, and the register shows the closure as
\code{manual}. A person may take it back by setting another state; the closing fields are
then cleared.

\smallskip
A closure made by a report is different. Setting another state on it is refused with 409
\code{CLOSED_BY_REPORT}: \enquote{A signed-off report closed this action --- to dispute it,
review the finding in that report}. The closure rests on findings of a signed-off report, so
the place to dispute it is the finding. Changing the finding there withdraws the report's
sign-off, and the same code that closed the action works the closure out again without that
report (\cref{sec:register:settle}). The console shows the status select of such an action
switched off, with a line under it that says where to dispute the closure.
\end{humanbox}

\subsection{Reading the register}

\route{GET}{/api/companies/{id}/actions} (\code{company_actions}) answers the register and the
count of actions in each shown state. The order is the one a person needs: the late actions
first, then the open and in-progress ones, then those called done, then the closed; within a
group by due date, then by the plan's priority. The list is cut at 500 rows; the counts are
not. \Cref{ch:console} describes the tab that shows it.


\section{Confirmation by later reports}\label{sec:register:settle}

This section is the centre of the register: how a later report confirms a finding, closes an
action or sends one back to work. One private function does all of it, \code{_settle(company_id,
now)}, and both \code{sign_off} and \code{withdraw} call it. It does not apply the newly signed
report to the register as a change. It works every confirmation out again from \emph{all} the
company's signed-off reports, each time. That is the property everything else follows from:
the confirmations and closures of the actions there are come out the same whatever order the
signatures were given in, and taking a signature back needs no undo log --- the function is simply run again without that report.

\subsection{What the function reads}

\code{_settle} takes the company's actions that have items and are either live or were closed
by a report. Actions a person closed by hand and actions without items are left alone. For the
obligations named by the items it loads every finding on them in the company's reports that are
finished and signed off, ordered by the time the report was \emph{made} (\code{created_at}),
and keeps for each obligation the list of verdicts: when the report was made, which report,
and the effective status.

For one action, the verdicts that count are those of reports made no earlier than the action's
own report (\code{since}; once that report is deleted, no earlier than the action itself). A report made before the action existed says nothing about work
done under it, even when it is signed off later.

\subsection{The rules}

For each item of the action, the function walks the obligation's verdicts in the order the
reports were made and looks for the report that \emph{first found the obligation settled with
no later report finding it open}: a settled verdict starts a run, an open verdict ends it.

\begin{enumerate}
  \item \textbf{A finding is confirmed} when the newest verdict on its obligation is settled.
    The item's \code{confirmed} is set to \code{{reportId, at}}, naming the first report of
    that unbroken run. A report that did not assess the obligation leaves no verdict and
    therefore confirms nothing and breaks nothing: a narrow report on a few documents in
    between does not disturb the register. A verdict of \status{unknown} is open --- a report
    that could not find evidence does not confirm.
  \item \textbf{An action is closed} when every one of its items is confirmed. Its state
    becomes \code{closed}, \code{closed_from} remembers the state it had, \code{closed_by}
    stays empty, and \code{confirmed_report} names the newest of the confirming reports ---
    the report that completed the set.
  \item \textbf{An action called done goes back to work} (\code{in_progress}) when the newest
    verdict on one of its obligations is open \emph{and} comes from a report made after the
    moment the owner called it done (\code{done_at}). The owner said the work was carried out;
    a report that looked afterwards still found the gap.
  \item \textbf{A done action is not sent back} by a report made before \code{done_at}, even
    when that report is signed off later: it never saw the finished work. It is not sent back
    by a report that did not assess the obligation either. An action that is \code{open} or
    \code{in_progress} is never moved by an open verdict; it simply stays.
  \item \textbf{A closed action stays closed.} For an action a report closed, only the
    verdicts up to the closing report are read, so a later report that finds the obligation
    open again does not reopen it. That report's own plan brings a \emph{new} action for the
    gap, because the obligation is no longer under a live action
    (\cref{sec:register:plan}), and, when it finds the obligation \status{missing} or
    \status{partial}, the report counts the gap as one that went back
    (\cref{sec:register:history}). The old action remains a true record: the gap was closed
    then, and reopened later.
  \item \textbf{A closure falls only with the evidence under it.} While the closing report is
    signed off, the items are read against the reports up to it; once it is not, against every
    report that still is. If every item is still confirmed, the action stays closed and
    \code{confirmed_report} names the newest confirming report that remains. If one is not,
    the closure does not stand: the action returns to the state it was closed from
    (\code{closed_from}), \code{confirmed_report} and \code{closed_at} are cleared, and the
    items are read once more against every signed-off report.
\end{enumerate}

The function answers four numbers --- findings newly confirmed, actions closed, actions sent
back to work, actions touched at all. A sign-off reports the first three as \code{confirmed},
\code{closed} and \code{reopened}; a withdrawal reports the fourth as \code{withdrawn}.

\subsection{A worked timeline}

\Cref{fig:register:timeline} follows one action through three reports. It is an example
constructed for this chapter from the rules above, not a record of real data; the test
\code{test_a_later_signed_report_is_the_re_test} in \file{tests/test_register.py} drives the
same kind of sequence through the API.

\begin{figure}[H]
\centering
\begin{tikzpicture}[x=1mm, y=1mm,
    cell/.style={align=center, text width=23mm, inner sep=1.5pt, font=\sffamily\scriptsize, text=body},
    rowlabel/.style={align=right, text width=19mm, anchor=east, inner sep=1.5pt,
                     font=\sffamily\scriptsize\bfseries, text=ink},
    ev/.style={man, text width=20.5mm, inner xsep=3pt, inner ysep=3pt, minimum height=11mm, font=\sffamily\scriptsize},
    st/.style={chk, text width=20.5mm, inner xsep=3pt, inner ysep=3pt, minimum height=7mm, font=\sffamily\scriptsize}]
  % columns
  \foreach \x/\n in {33.5/1, 58.5/2, 83.5/3, 108.5/4, 133.5/5} {\coordinate (c\n) at (\x,0);}
  % time arrow
  \draw[flow] (21,0) -- (147,0);
  \node[rowlabel] at (20,0) {In this order};
  \foreach \n in {1,3,4} {\fill[ink] (c\n) circle[radius=.9];}
  \node[note, anchor=south] at (33.5,1) {report 1 made};
  \node[note, anchor=south] at (83.5,1) {report 2 made};
  \node[note, anchor=south] at (108.5,1) {report 3 made};
  % events
  \node[rowlabel] at (20,-10) {A person};
  \node[ev] at (33.5,-10) {signs report 1 off};
  \node[ev] at (58.5,-10) {calls the action done};
  \node[ev] at (83.5,-10) {signs report 2 off};
  \node[ev] at (108.5,-10) {signs report 3 off};
  \node[ev] at (133.5,-10) {withdraws report 3's sign-off};
  % obligation X
  \node[rowlabel] at (20,-24) {Obligation X};
  \node[cell] at (33.5,-24) {\status{missing}\\ item of the action};
  \node[cell] at (58.5,-24) {---};
  \node[cell] at (83.5,-24) {\status{met}\\ confirmed by report~2};
  \node[cell] at (108.5,-24) {\status{met}\\ still confirmed by report~2};
  \node[cell] at (133.5,-24) {confirmed by report~2};
  % obligation Y
  \node[rowlabel] at (20,-38) {Obligation Y};
  \node[cell] at (33.5,-38) {\status{partial}\\ item of the action};
  \node[cell] at (58.5,-38) {---};
  \node[cell] at (83.5,-38) {\status{partial}\\ not confirmed};
  \node[cell] at (108.5,-38) {\status{met}\\ confirmed by report~3};
  \node[cell] at (133.5,-38) {confirmation gone};
  % state
  \node[rowlabel] at (20,-52) {The action};
  \node[st] (s1) at (33.5,-52) {\code{open}};
  \node[ev, minimum height=7mm] (s2) at (58.5,-52) {\code{done}};
  \node[st] (s3) at (83.5,-52) {\code{in_progress}};
  \node[st] (s4) at (108.5,-52) {\code{closed}};
  \node[st] (s5) at (133.5,-52) {\code{in_progress}};
  \draw[flow] (s1) -- (s2);  \draw[flow] (s2) -- (s3);  \draw[flow] (s3) -- (s4);  \draw[flow] (s4) -- (s5);
  \node[cell, text=muted] at (33.5,-61) {the plan comes in; the due date counts from here};
  \node[cell, text=muted] at (58.5,-61) {\code{done_at} is stamped};
  \node[cell, text=muted] at (83.5,-61) {sent back: Y is open in a report made after \code{done_at}};
  \node[cell, text=muted] at (108.5,-61) {both items confirmed; closed by report~3};
  \node[cell, text=muted] at (133.5,-61) {back to the state it was closed from};
  % separators
  \begin{scope}[on background layer]
    \foreach \x in {46, 71, 96, 121} {\draw[draw=rule, line width=.4pt] (\x,-4) -- (\x,-66);}
    \foreach \y in {-17, -31, -45} {\draw[draw=rule, line width=.4pt] (0,\y) -- (146,\y);}
  \end{scope}
\end{tikzpicture}
\actorlegend
\caption{One action, two obligations, three reports. Each column is one act of a person;
the rows show what the newest signed-off report says about each obligation, what the
register has confirmed, and the state of the action. Every light blue box is worked out by
code at a sign-off or at its withdrawal.}\label{fig:register:timeline}
\end{figure}

Read column by column:

\begin{enumerate}
  \item \textbf{Report 1 is signed off.} Its plan holds an action for two findings: obligation
    X is \status{missing}, obligation Y is \status{partial}. The action comes in as
    \code{open} with two items.
  \item \textbf{The owner calls it done.} \code{done_at} is stamped. Nothing is confirmed:
    this is a statement.
  \item \textbf{Report 2, made after that, is signed off.} It finds X \status{met} and Y still
    \status{partial}. The item on X is confirmed by report~2. The item on Y is not, and because
    report~2 was made after \code{done_at}, the action goes back to \code{in_progress}. The
    sign-off answers \code{confirmed: 1}, \code{reopened: 1}, and, when report~2's plan names its
    own open finding on Y again, counts that finding as \code{covered}. Had report~2 not assessed Y at all, the action would have
    stayed \code{done}.
  \item \textbf{Report 3 is signed off.} It finds both met. Y is confirmed by report~3; X stays
    confirmed by report~2, the first report of its unbroken run. Every item is confirmed, so
    the action is closed by report~3, from the state \code{in_progress}.
  \item \textbf{Report 3's sign-off is withdrawn} --- by hand, or because a reviewer changed
    one of its findings. Without report~3, Y has no settled verdict: its confirmation is
    cleared and the action returns to \code{in_progress}. X stays confirmed, because report~2
    is still signed off. Signing report~3 off again closes the action again.
\end{enumerate}

\subsection{The order the reports were made in, not the order they were signed in}

Reports are not always signed off in the order they were made. Because \code{_settle} sorts
verdicts by the report's creation time and reads them all, an old report signed late cannot
overrule a newer one. The test \code{test_the_reports_speak_in_the_order_they_were_made}
signs a fourth report, which finds everything met, before the third, which still finds one
obligation open: the fourth closes the action, and the third, signed afterwards, changes
nothing and adds no action for the finding the newer report finds met. The same test then
deletes reports: with the third deleted the closure stands, because the fourth found the
same; with the fourth deleted as well, the evidence under the closure is gone and so is the
closure.

\subsection{What a withdrawal does not undo}

The actions a report brought in stay when its sign-off is withdrawn or the report is deleted:
in the words of the code, \enquote{they are the company's now}. Only confirmations and the
closures resting on them follow the evidence. An action whose report was deleted keeps its
items and loses only the link to the report. Signing the report off again brings nothing
twice.


\section{Identity: an obligation is known by its id}\label{sec:register:identity}

Everything in \cref{sec:register:settle} joins a finding of one report to a finding of
another through \code{obligationId}. The join is exact, with no matching by title or by
provision; it works only while an obligation keeps its id from one report to the next. When
that holds is decided outside the register, by how obligations are stored.

\begin{itemize}
  \item \textbf{An obligation of an act or a standard} keeps its id across a re-reading of its
    act once a person has approved it. Re-extraction deletes and replaces only the rows still
    \code{proposed} (\file{corpus/obligations.py}); approved and rejected rows are never
    touched. A proposed obligation that is extracted again is a new row with a new id.
    \Cref{ch:obligations} describes the extraction and the review.
  \item \textbf{A company's own requirement} keeps its id when its policy is read again and
    the passage that states it is unchanged (\code{read_requirements} in
    \file{analysis/internal.py}): the new row takes the id of the proposed row with the same
    quoted passage. A passage that was edited yields a new requirement.
\end{itemize}

When the obligation of an unconfirmed item no longer exists, no later report will ever carry
its id, and the action could wait for a confirmation for ever. The register says so instead.
\code{standing(actions)} looks up which of the items' obligations still exist, and
\code{serialize_company_action} serves an unconfirmed item whose obligation is gone with
\code{obligationId: null}. The action then carries \code{byHandOnly: true}, as does every
action that has no items at all. The console prints one quiet line on such an action ---
\enquote{No report can confirm this action --- only a person closes it, giving the reason in
the note} --- and, on the item, why: its act, or for a requirement of the company's own its
policy, was read again since. The stored item keeps the id; only the answer hides it.

This is why the review of the register of obligations comes before relying on confirmations,
and why an act whose proposed obligations are under open actions should not be extracted
again before they are reviewed (\cref{ch:ops}). On 13~November 2026 the live register of obligations held 5,960
obligations, 5,938 of them approved and 12 still proposed; \cref{ch:obligations} states how
that review was done.


\section{History and movement}\label{sec:register:history}

Two smaller functions use the same join to look backwards from a report. Unlike the
confirmations, they read every \emph{finished} earlier report of the company, signed off or
not: they describe what the reports said, not what a person accepted.

\paragraph{A finding's history.} \code{history(report)} answers, for each obligation of the
report's findings, the last times an earlier finished report assessed it, the newest first, at
most \code{HISTORY}~=~4 entries of \code{{reportId, date, status}} with the effective status.
A report that did not assess the obligation leaves no entry. \code{serialize_report} attaches
the list to every finding of a full report as \code{history}; it is worked out when the report
is served, so it reflects later overrides in the earlier reports. The console shows the first
three entries on the finding's card as a row of earlier verdicts, each a link to that report.

\paragraph{What moved.} \code{changes(report)} compares the report with the company's earlier
finished reports, obligation by obligation, and answers three lists, each as a count with up
to five titles. Here a \emph{gap} is a finding that is \status{missing} or \status{partial};
\status{unknown} is neither a gap nor settled.

\begin{table}[H]
\caption{How \code{changes} classifies an obligation of the report}\label{tab:register:changes}
\begin{tabularx}{\textwidth}{@{}L{20mm}L{44mm}Y@{}}
\toprule
\hd{Key} & \hd{Earlier verdict → now} & \hd{Reads as} \\
\midrule
\code{closed} & a gap → settled & A gap that is met now. \\
\code{back} & settled → a gap & An obligation that was met and is a gap again. \\
\code{new} & not a gap and not settled, or never assessed → a gap & A gap on an obligation
  that no earlier report assessed, or that was \status{unknown} before. \\
\bottomrule
\end{tabularx}
\end{table}

The earlier verdict is the one of \emph{the last earlier report that assessed that
obligation}, not of the report just before. The reason is practical: a report may be run on a
few documents or a few acts. If every obligation were compared with the report just before,
a narrow report in between would make everything it did not assess look new in the next full
report. The test \code{test_a_narrow_report_in_between_does_not_make_the_rest_look_new} fixes
this behaviour. The answer still names the report just before (\code{since}, \code{sinceDate})
as the date to speak of. Findings on the company's own requirements are left out of the
movement, as they are counted apart everywhere, and the function answers nothing when the
company has no earlier finished report.

The result is used in three places, all owned by \cref{ch:report}. It is stored in the
report's statistics as \code{stats.changes} when the report finishes (\code{report_stats} in
\file{analysis/pipeline.py}). It is added to the input of the summary step as a block headed
\texttt{SINCE THE PREVIOUS REPORT}, with the three counts and their titles, and the summary
prompt is told to close its first paragraph, the overall assessment, with one sentence on the
movement. And the
report's PDF prints one line under its results at a glance (\file{analysis/pdf.py},
\code{_glance}): \enquote{Since the previous report (date): gaps closed -- n, new gaps -- n,
went back -- n.} In all three the numbers are computed by code; the model only words the
sentence.


\section{Changes in the law}\label{sec:register:law}

A report is an assessment against the versions of the acts that were current when it ran.
When an act changes afterwards, some of the report's findings rest on provisions that no
longer read as they did. The ingestion records what changed, and the register's
\code{law_changes} names those findings.

\paragraph{What is recorded.} When the ingestion job activates a new version of an act, it
stores in \code{legal_document_versions.changes} what the version changed against the one it
replaces (\code{version_changes} in \file{corpus/ingest.py}): \code{{from, changed, added,
removed}}, where the three lists hold provision paths in the act's order. The comparison is
article by article --- for an act without articles, by its top-level points --- on the text
with white space normalised: \code{changed} are paths present in both versions with different
text, \code{added} those only in the new one, \code{removed} those only in the old one. An
act's first version has no record (null); nor has a version that is processed again in place,
which would be compared with itself. No model reads the two versions; it is a comparison
of stored text. \Cref{ch:corpus} describes ingestion and versions.

\paragraph{What the route lists.} \route{GET}{/api/companies/{id}/law-changes}
(\code{law_changes}) takes the company's latest finished report and, for every act that report
assessed (as \code{assessed} or \code{upcoming}) whose current version is no longer the one the
report read, merges the records of all versions since. Each act is answered with the version
number the report read and the current one with its date, the three lists, and
\code{findings}: the report's findings under that act whose citations rest on a changed or
removed provision, each with its id, title, effective status and the cited path. A citation
rests on an article when its path is the article's path or continues it without a further
digit, so \texttt{Art.\ 68(1)} rests on \texttt{Art.\ 68} and \texttt{Art.\ 680} does not. Added
provisions touch no existing finding and are listed for information.

The flag \code{recorded} is false when a version in between has no record. The act then
changed and the system cannot say where, and it says exactly that rather than show empty lists
that would read as \enquote{nothing changed}: the console prints \enquote{The act changed, but
which provisions was not recorded}, and the chat tool is told that the whole act is to be
looked at again.

The console shows the list under the register, in the section \ui{Teisės aktų pokyčiai}
(changes in the law), with each finding linked to its card in the report. The list changes
nothing: it does not reopen an action or alter a finding. It names where a person should look,
and a new analysis assesses the company against the text now in force.

On 13~November 2026 no stored version on the live system carried a record of changes: no new
version of an act had been activated since the recording was introduced. The mechanism is
exercised by the test \code{test_a_new_version_of_an_act_names_the_findings_to_look_at_again},
which rewords one article, drops one and adds one, and checks the three lists and the findings
named; it has not yet met a real amendment.


\section{Reading the register: the chat and the Excel export}\label{sec:register:reading}

\subsection{The chat assistant reads and never writes}

The chat assistant has one tool for the register, \code{getActionRegister}, in the assistant
service's \file{llm/tools/actions.js}. It calls the two reading routes --- the register, then
the changes in the law --- and returns a compact answer: the counts by shown state; up to
\code{MAX_ACTIONS}~=~20 actions, the closed ones only when that state is asked for, each with
its state, owner, due date and up to \code{MAX_ITEMS}~=~6 of the findings it closes;
\code{closedOnlyByHand} where the API says \code{byHandOnly}; and, when acts moved, the changes
in the law with the findings to re-check. If the law-changes route fails, that part is left
out and the register is still answered.

The tool takes no argument that could change anything, and no other chat tool signs a report
off or edits an action; a test in \file{tests/tools-actions.test.mjs} asserts that the only
tool whose name mentions an action, a sign-off or a register is this one. What the model may
\emph{say} is constrained twice: by the tool's description and its \code{next} hint, returned
with every answer, and by a rule of the chat prompt.

\begin{promptexcerpt}{chat.assistant}{rule 13, on the register}
[…]
The register moves only when a person signs a report off on the report page: its plan comes in, and a later signed-off report that finds an action's obligations met closes it. You cannot open, close or edit an action, and you never call one closed because a document was uploaded or the user says it is done – "done" there means carried out and awaiting that re-test; tell the user to sign the next report off, or to close the action by hand with a reason on the company's Actions tab.
[…]
\end{promptexcerpt}

The rule exists because the tempting answer is the wrong one. A user who has just uploaded
the missing procedure and asks \enquote{is it closed now?} would be pleased by a yes, and a
model inclined to be helpful could give it. The register's answer is that the upload is not
evidence until a report has assessed it and a person has signed that report. \Cref{ch:chat}
describes the chat's tools and prompt as a whole.

One indirect path deserves stating. The chat can record a reviewer's verdict on a finding, on
the user's instruction and with the user's reason (the tool \code{overrideFinding},
\cref{ch:chat}). It sends the same request as the console's review dialog, so on a signed-off
report it has the same consequence: the sign-off is withdrawn, and what it confirmed in the
register goes back. The chat still cannot give a sign-off.

\subsection{The Excel export}

\route{GET}{/api/companies/{id}/actions/export} answers the whole register as a workbook
(\code{render_actions_xlsx} in \file{analysis/exports.py}), in Lithuanian or English by the
\code{language} parameter, laid out in the columns of the course's findings-and-actions
register. One row is one action, in the order of the list: a running number; the action with
its description; the findings it is to close, one line each with the act and the provision
and the word \enquote{confirmed} once a later report confirmed it; the owner; the due date;
the state as shown, late included; the priority; the date of the report it came from; the
date it was closed; how it was closed (\enquote{confirmed by a later report} or
\enquote{closed by a person}); and the note. Every text cell is stored as text, never as a
formula. The download is written to the audit log as \code{actions.export}. The chat tool
words what an action closes the same way as the workbook, so the two agree.


\section{Where the line is drawn}\label{sec:register:line}

\begin{codebox}{No model is called anywhere in the register}
\file{analysis/register.py} imports the report and corpus models, the database transaction and
a clock. It does not import the function that calls a model, and nothing it calls does.
Bringing a plan in, confirming a finding, closing an action, sending one back to work,
deriving \code{late}, a finding's history, the movement between reports and the changes in
the law are all comparisons of stored ids, statuses, dates and texts. Given the same reports,
the same signatures and the same edits by people, in the same order, the register always comes
out the same.

\smallskip
Its \emph{inputs} are another matter, and should be named. The titles, the owner role, the
timeframe and the category of an action were written by a model in the report's plan, and
the status of a finding is a model's verdict unless a reviewer overrode it. The register adds
no judgement of its own to them: it only records which of them a person signed.
\end{codebox}

\begin{humanbox}{What only people do here}
A person signs a report off and may take the signature back. A person owns an action, sets
its due date and says where it stands. A person may close an action by hand, with a reason
that stays on the record, and may reopen what a person closed. A person disputes a report's
closure at the finding it rests on. Nothing in the register happens on the system's own
initiative: every change follows from one of these acts.
\end{humanbox}

The line stands here for three reasons. The first is accountability: an action register says
what a company committed itself to and whether it kept the commitment, and both statements
should trace back to a named person and a dated report, not to a model's wording on one day.
The second is reproducibility: a model asked twice whether an action is closed may answer
twice differently, while a join on an obligation's id cannot. The third is that the judgement
a model \emph{is} good for has already been made, checked and reviewed elsewhere --- in the
finding (\cref{ch:report}). The register reuses that judgement and adds none, so it adds no
new place where a model could be wrong. \Cref{ch:control} gathers the system's controls and
places this one among them.


\section{What was tried, what is tested, and a known limitation}\label{sec:register:tried}

\begin{lessonbox}{Tried on a copy of the database, 12 November 2026}
The register was run against the real reports of the sample company, \demo, on a copy of the
database, with no model call. Signing the company's full report off brought its 25 plan
actions in with the 587 open findings they are to close, every one with its obligation's id.
Signing the later small report off added 7 actions and left 6 out as already under way. Signed the other way round --- the newer report first --- the register held
the same 653 findings, none under two actions.

\smallskip
A re-test was then staged: the full report was copied and some of its findings set to met.
Signing the copy off confirmed 7 findings, closed the action whose findings were all met, and
sent back to work the one action that had been called done and still had a finding open.
Withdrawing that sign-off undid exactly that. Deleting the copy after a second re-test left
the closures standing on the second. These are observations of one day on one company; the
re-test was staged by editing findings, not produced by a second real analysis after real
remediation.
\end{lessonbox}

\paragraph{Tests.} \file{tests/test_register.py} holds seven tests that drive the register
through the API's routes with fake models: a sign-off brings the plan in once; a later signed
report is the re-test; the reports speak in the order they were made; a person owns, dates,
moves and closes an action, and the export; a finding's history and a report's movement; a
narrow report in between; and a new version of an act. The authorisation matrix
(\file{tests/test_authz_matrix.py}) covers the register's routes, and the assistant service
tests the chat tool. \Cref{ch:testing} describes the suites.

\paragraph{A known limitation.} \code{standing()} asks only whether an obligation's row still
exists. A report, however, assesses the obligations of an act's current version that are not
rejected. An obligation that an administrator rejects after it came under an action therefore
still exists, is never assessed again, and is not marked \code{byHandOnly}: its action waits
for a confirmation that cannot come, until a person notices and closes it by hand. On
13~November 2026 none of the ten rejected obligations was under an action. The same holds, by
the same reading of the code, for an approved obligation that a new version of its act left
behind because the provision it cites changed: \code{carry_forward} leaves such a row on the
old version, where no later report reads it. In that second case, until the company's next
report, the act has a newer version than its latest report read, so the change is at least
listed under the changes in the law (\cref{sec:register:law}), where a person is pointed to it.
